% =====================================================================
%  БАЗОВЫЙ ШАБЛОН РЕЗЮМЕ (RU)
%
%  1. Скопировать в resumes/<семейство>_ru.tex
%  2. Заполнить блок «ДАННЫЕ» значениями из profile.json (комплект контактов —
%     строго под рынок; комплекты не смешивать)
%  3. Наполнить секции фактами из facts.md. КАЖДЫЙ содержательный блок несёт
%     комментарий % FACT: <id> — по нему ats_check.py проверяет провенанс.
%
%  Ничего сверх facts.md не добавлять. Формат намеренно простой: одна колонка,
%  без таблиц, колонок, иконок и фото — иначе строгие парсеры (Workday) ломаются.
% =====================================================================
\documentclass[a4paper,10pt]{article}

\usepackage{cmap}
\usepackage[T2A]{fontenc}
\usepackage[utf8]{inputenc}
\usepackage[russian,english]{babel}
\usepackage{geometry}
\usepackage{enumitem}
\usepackage{hyperref}
\usepackage{xurl}
\usepackage{microtype}
\usepackage{titlesec}

\geometry{left=1.35cm,right=1.35cm,top=1.25cm,bottom=1.25cm}
\setlength{\parindent}{0pt}
\setlist[itemize]{leftmargin=1.15em, itemsep=2pt, topsep=2pt, parsep=0pt}

\titleformat{\section}{\large\bfseries}{}{0em}{}[\titlerule]
\titlespacing{\section}{0pt}{6pt}{4pt}

\hypersetup{colorlinks=true, urlcolor=black, linkcolor=black}
\Urlmuskip=0mu plus 1mu
\emergencystretch=2em
\sloppy

% ОБЯЗАТЕЛЬНО: без этого кириллица извлекается из PDF кракозябрами
% и ATS видит мусор вместо текста.
\input{glyphtounicode}
\pdfgentounicode=1

% Часто повторяющиеся сокращения. \mbox защищает от переноса внутри токена:
% перенос после «/» даёт при извлечении битые токены вида «C/ C++».
% Не заменять на \allowbreak.
% \newcommand{\CCPP}{\mbox{C/C++}}

% ---------- ДАННЫЕ (из profile.json) ----------
\newcommand{\CVName}{Имя Фамилия}
\newcommand{\CVTitle}{Заголовок из facts.md §1}
\newcommand{\CVEmail}{name@example.com}
\newcommand{\CVPhone}{+7 000 000 00 00}
\newcommand{\CVTelegram}{example}

\begin{document}

% ---------- ШАПКА ----------
% FACT: IDENT-name CONTACT-RF TITLE-XXX
{\LARGE \textbf{\CVName}}\\
\textbf{\CVTitle}\\
Email: \href{mailto:\CVEmail}{\CVEmail} \quad
Телефон: \CVPhone \quad
Telegram: \href{https://t.me/\CVTelegram}{@\CVTelegram}

% ---------- О СЕБЕ ----------
% Абзац SUM-... из facts.md §2. Верхняя треть резюме — самая весомая для
% ранжирования: ключевые слова вакансии размещать здесь.
\section*{О себе}
% FACT: SUM-XXX

% ---------- НАВЫКИ ----------
% Группы из facts.md §7. Только verified/basic. [gap]-инструменты не писать.
% Каждый заявленный навык должен подтверждаться буллетом в «Опыт» или «Проекты».
\section*{Навыки}
% FACT: SKL-XXX
% \textbf{Языки:} ...\\
% \textbf{Инфраструктура:} ...

% ---------- ОПЫТ ----------
% Из facts.md §4. Порядок «организация -> должность -> даты» парсится корректно
% строгими ATS. Даты писать «месяц год», голых годов избегать.
\section*{Опыт}
% FACT: EXP-XXX
% \textbf{Организация} \hfill \textit{март 2023 -- н.в.}\\
% \textit{Должность}
% \begin{itemize}
%   \item ...
% \end{itemize}

% ---------- ПРОЕКТЫ ----------
% Из facts.md §5.
\section*{Проекты}
% FACT: PRJ-XXX
% \textbf{Название} \hfill \textit{стек / контекст}
% \begin{itemize}
%   \item ...
% \end{itemize}

% ---------- ОБРАЗОВАНИЕ ----------
% Из facts.md §3.
\section*{Образование}
% FACT: EDU-XXX
% \textbf{Учреждение} \hfill \textit{2017 -- 2021}\\
% Программа

% ---------- ДОСТИЖЕНИЯ ----------
% Из facts.md §8, отфильтровать по тегам семейства. Секцию можно убрать,
% если достижений нет или они нерелевантны.
\section*{Достижения}
% FACT: AWD-XXX
% \begin{itemize}
%   \item ...
% \end{itemize}

\end{document}
